\section{Finite whole canonical values after the actual first NC step}
\label{sec:hironaka-normal-crossings-first-step-finite-values}

Assume $n>0$. Retain the unchanged original NC marking-one datum $M_H$,
its positive highest branch count $r$, and a genuine blowup of its entire
highest stratum. The invariant throughout is exactly
$R_n=\mathrm{canonicalInvariant}(n)$. No equality of invariant depths,
compactness, prescribed finite value image or later history is assumed.

\begin{lemma}[The actual coordinate union has its intrinsic coordinate count]
\label{hir:general-nc-coordinate-carrier-intrinsic-count}
\leanok
\lean{Hironaka.HironakaNormalCrossingsFirstStepFiniteValues.coordinateCarrier_branchMultiplicity}
For the actual coordinate union $H_S$, the intrinsic branch count at $w$
equals the number of selected coordinates vanishing at $w$.
\uses{hir:general-nc-standard-coordinate-carrier}
\end{lemma}
\begin{proof}
Translate $w$ to zero and restrict to the actual open neighbourhood on
which every selected coordinate nonzero at $w$ remains nonzero. The whole
carrier on this neighbourhood is then the union of precisely the selected
coordinates that vanish at $w$. Translation leaves those coordinates
unchanged because their values at $w$ are zero. This gives a genuine
centred NC presentation with exactly that finite index set. The intrinsic
branch-count theorem for an actual NC presentation identifies its
cardinality with the intrinsic count.
\end{proof}



\begin{lemma}[The fixed entire coordinate-model first history is genuinely canonical]
\label{hir:general-nc-entire-coordinate-first-history-canonical}
\leanok
\lean{Hironaka.HironakaNormalCrossingsFirstStepFiniteValues.coordinate_model_first_isCanonical}
For the union of the last $r$ coordinate hyperplanes, with $r>0$ and
exactly $r$ selected indices, the global highest stratum is the whole
coordinate centre. Every genuine blowup of that centre is an actual
first canonical test step for the unchanged original coordinate datum.
\uses{hir:general-nc-coordinate-carrier-intrinsic-count,hir:general-nc-ambient-canonical-initial-full-maximum}
\end{lemma}
\begin{proof}
The preceding formula bounds every coordinate count by $r$ and gives
count $r$ at zero. A point has count at least $r$ precisely when all
selected coordinates vanish: its selected vanishing-coordinate subset
has full cardinality and hence equals the selected set. Thus its whole
highest stratum is exactly the coordinate centre. Apply the actual
ambient-dimensional initial full-maximum theorem to this original NC
carrier, using zero for attainment. The centre of the given genuine
blowup is consequently the entire whole maximum. The canonical-step
constructor gives the first canonical history, and proved maximum
admissibility gives its actual test-history property.
\end{proof}

\begin{theorem}[Finite whole model values on the actual exceptional fibre]
\label{hir:general-nc-entire-model-finite-whole-fibre-values}
\leanok
\lean{Hironaka.HironakaNormalCrossingsFirstStepFiniteValues.finite_coordinate_first_fibre_values}
The whole value image of $R_n$ on the actual transformed cosupport over
zero in that same fixed entire coordinate-model first blowup is finite.
\uses{hir:general-nc-entire-coordinate-first-history-canonical,hir:general-nc-finite-geometric-fibre-orbit-representatives}
\end{theorem}
\begin{proof}
The preceding lemma supplies genuine entire-model canonicity. Actual
diagonal base maps preserve the whole original coordinate carrier and
therefore its entire vanishing-ideal marked record, including the empty
boundary and marking one. Their actual lifted one-step history
isomorphisms preserve the literal transformed cosupport by the whole
transform law. The invariant's genuine tower-isomorphism law therefore
makes the entire first-step values equal at each point and its actual
diagonal image. The finite geometric representative producer supplies
one fixed finite set $Q$ and, for every point of the fibre, such an actual
lift taking it into $Q$. Its value thus belongs to the image of $Q$ under
the fixed entire-model value function. This image is finite. Unused
indexed points in $Q$ cause no difficulty, since only containment in
that finite image is needed.
\end{proof}

\begin{lemma}[The actual original first history compares with the same entire model]
\label{hir:general-nc-original-first-value-entire-coordinate-model}
\leanok
\lean{Hironaka.HironakaNormalCrossingsFirstStepFiniteValues.exists_highest_chart_entire_canonical_value}
At a point over the original highest stratum, the original whole
first-step value equals the fixed entire coordinate-model value at an
actual point over zero, with transformed cosupport retained.
\uses{hir:general-nc-first-step-whole-canonical-chart-value,hir:general-nc-entire-coordinate-first-history-canonical}
\end{lemma}
\begin{proof}
Choose the highest-stratum chart at the actual image of the blown-up
point. The actual one-step chart-history producer constructs the genuine
history isomorphism, whole original-record identity and whole value
equality on the literal coordinate open history. Its exact commuting
projection takes the actual point to a coordinate-model point over zero,
because the chosen chart is centred at its original image. The whole
history transform and cosupport restriction identities put this actual
image point in the literal transformed coordinate cosupport.

The fixed entire-model history is canonical by the preceding lemma.
Apply its proved open restriction law to the coordinate open target.
Actual vanishing-ideal restriction identifies the pulled-back original
record with the literal original coordinate-open datum. The whole
cosupport restriction identity supplies its point hypothesis. Composing
this restriction equality with the actual chart-history value equality
gives the asserted equality with the SAME entire-model history. Only
after this step is diagonal normalization used; no scaling is required
to preserve the original chart's open target.
\end{proof}

\begin{lemma}[Finite original initial values and the actual off-centre branch]
\label{hir:general-nc-initial-and-first-off-centre-finite-values}
\leanok
\lean{Hironaka.HironakaNormalCrossingsFirstStepFiniteValues.finite_initialValue_image,Hironaka.HironakaNormalCrossingsFirstStepFiniteValues.finite_first_off_centre_values}
Every proved invariant with dimension bound at least $n$ has a finite
whole initial value image on an actual original NC carrier. For $R_n$,
the actual first-step value image off the first centre is finite.
\uses{hir:general-nc-generic-initial-value-flat-chart,hir:general-nc-ambient-canonical-initial-full-maximum}
\end{lemma}
\begin{proof}
At every original carrier point, its actual NC presentation selects a
subset $S$ of the finite coordinate index type. Generic whole initial
chart transport identifies its entire value with that of $H_S$ at zero.
There are only finitely many such selected subsets, so the original
value image is contained in a finite range of actual model values.
For the actual first canonical blowup, the proved off-centre invariant
law leaves every whole value unchanged at its original image. Actual
test-transform cosupport projection puts that image in the unchanged
original cosupport, which equals the original carrier. Thus all
off-centre first-step values belong to the same finite initial image.
\end{proof}

\begin{theorem}[Finite whole values after the actual first original NC blowup]
\label{hir:general-nc-first-step-finite-entire-whole-values}
\leanok
\lean{Hironaka.HironakaNormalCrossingsFirstStepFiniteValues.finite_first_canonical_value_image,Hironaka.HironakaNormalCrossingsFirstStepFiniteValues.exists_first_canonical_whole_maximum}
From the actual original NC carrier, its positive attained highest
stratum and a genuine blowup of that stratum, the entire first transformed
cosupport has a finite whole $R_n$ value image. In particular a nonempty
first transformed cosupport has an internally attained whole maximum.
\uses{hir:general-nc-entire-model-finite-whole-fibre-values,hir:general-nc-original-first-value-entire-coordinate-model,hir:general-nc-initial-and-first-off-centre-finite-values}
\end{theorem}
\begin{proof}
An actual highest-stratum NC chart at an attained point has precisely
the last $r$ branches. The intrinsic chart-count theorem therefore
derives that this selected coordinate set has cardinality $r$ from the
original geometry. The actual highest-stratum blowup producer for the
corresponding whole coordinate NC carrier internally chooses a fixed
genuine entire-model blowup. Its actual whole highest centre is the
coordinate centre by the coordinate-count lemma.

Split the original first transformed cosupport into points over the
centre and points off it. In the first case, centre the actual original
chart at the point's original image. Whole history comparison with the
SAME fixed entire model puts its value in that model's finite whole
fibre image. In the second case, the preceding off-centre lemma gives
a finite whole value image. Their union is finite, so the whole original
first-step image is finite. If the transformed cosupport is nonempty,
so is this finite image. Its finite linear order has an actual greatest
element; choosing its actual preimage yields a transformed-cosupport
point attaining the whole maximum. No compactness or maximum certificate
is used, and no third or later step or terminal resolution follows
from this finite first-stage statement alone.
\end{proof}

\begin{theorem}[Actual first-prefix resolution or a genuine second full canonical blowup]
\label{hir:general-nc-actual-second-full-canonical-step}
\leanok
\lean{Hironaka.HironakaNormalCrossingsFirstStepFiniteValues.first_canonical_step_terminal_or_second}
For the same original datum and actual highest-stratum first blowup,
either that actual first tower is a marked test resolution, or its
nonempty whole full maximum internally supplies an actual second
canonical test blowup. The original tower, chronological boundary and
unchanged original invariant are retained literally.
\uses{hir:general-nc-first-step-finite-entire-whole-values}
\end{theorem}
\begin{proof}
The initial whole-maximum theorem makes the actual original first tower
canonical and hence a test tower. If its transformed cosupport is empty,
these exact two properties constitute a test resolution. Otherwise the
preceding finite-image theorem supplies a point attaining its whole
canonical maximum, so that entire maximum locus is nonempty. The proved
invariant admissibility law makes this same entire maximum a closed
smooth admissible centre for the unchanged current marked datum.
Cosupports along actual test towers of a genuine original invariant
domain have empty interior; hence their whole maximum loci have empty
interior as well. Actual smooth-centre blowup existence therefore
constructs a genuine blowup of that same entire maximum. Appending it
to the literal original first tower gives an actual test step by the
derived admissibility and a canonical step by the derived full-maximum
identity. No second centre, value, history or existence certificate is
input. This produces an actual second step when needed, without
asserting a third step or global termination.
\end{proof}
